\begin{workedexample}[Worked Example: Shannon capacity]
A channel has bandwidth $B = 20\ \text{MHz}$ and $\text{SNR} = 100$
($20\ \text{dB}$). Its capacity is
\[ C = B\log_2(1+\text{SNR}) = 20\times10^{6}\,\log_2(101) = 20\times10^{6}\times6.66 = 133\ \text{Mbit/s}. \]
A practical link runs a few dB below this, choosing a modulation and code whose
required $E_b/N_0$ fits the available SNR.
\end{workedexample}

\begin{keytakeaways}
\begin{itemize}[leftmargin=1.4em,itemsep=2pt]
\item $\text{SNR} = (E_b/N_0)(R_b/B)$; BER depends on $E_b/N_0$, e.g.\ QPSK $= Q(\sqrt{2E_b/N_0})$.
\item Shannon capacity $C = B\log_2(1+\text{SNR})$ is the unbeatable ceiling.
\item Pulse shaping fits the spectral mask; higher-order modulation trades SNR for spectral efficiency.
\end{itemize}
\end{keytakeaways}

\begin{exercises}
\item Shannon capacity of a $5\ \text{MHz}$ channel at $\text{SNR} = 31$ ($15\ \text{dB}$)?
\item Does BER depend on absolute SNR or on $E_b/N_0$?
\item Why can't symbols be made arbitrarily sharp in time?
\end{exercises}

\furtherreading{Proakis~\cite{proakis}; Sklar~\cite{sklar}.}
